\documentclass[11pt]{article}

\usepackage[T1]{fontenc}
\usepackage{lmodern}
\usepackage[margin=1in]{geometry}
\usepackage{amsmath}
\usepackage{booktabs}
\usepackage{longtable}
\usepackage{array}

\title{Dataset Features in the EPSS-TPG Experiments}
\author{}
\date{}

\begin{document}
\maketitle

\section{Overview}

The experiments use two families of datasets. The first is the
Sec4AI4Aec-EPSS-Enhanced LLM-summary corpus: GPT, Gemma, Mistral and
DeepSeek (refreshed in this round, with Mistral as a new addition),
plus the older NVD/KEV labelled corpus where the binary label is
membership in the CISA Known Exploited Vulnerabilities (KEV) catalog.
The second family is the Megavul branch: GPT, Mistral and Gemma
generated against the Megavul commit-level vulnerability corpus.
Megavul rows are commit-grounded: each row carries the actual
vulnerable code, the patched code, and the commit URL alongside the
LLM summaries. The Megavul datasets share a common 29-column schema
and a different label inventory from the social-media datasets
(Section~\ref{sec:megavul}).

The pipeline does not consume every raw column directly. Each dataset
is converted into a unified labelled-records form, and the training
code then builds two kinds of input from those records:

\begin{itemize}
    \item a Text Property Graph (TPG), built from the CVE description
    and, optionally, one or more LLM summary columns;
    \item a tabular feature vector, used only when the hybrid model is
    enabled.
\end{itemize}

For EPSS soft-label prediction the target is EPSS itself, so the EPSS
score and percentile are removed from the tabular input via the
\verb|--no-epss-feature| flag. The NVD/KEV binary experiments target
KEV membership instead, so EPSS can stay in the tabular branch
without leaking the label.

\section{Datasets}

The four social-media LLM-summary datasets share most of their schema.
They are not interchangeable: each dataset's summary columns were
generated by a different LLM, the missing-summary patterns differ by
dataset, and DeepSeek uses an older column naming.

\subsection{GPT Summary Dataset}

GPT contains $9{,}218$ records covering $9{,}218$ unique CVEs (one row
per CVE). Each row carries the NVD description, the CVSS components,
EPSS score and status, social and GitHub source counts, public-code
availability, and three GPT-generated summary columns.

\begin{table}[h]
\centering
\small
\begin{tabular}{ll}
\toprule
Property & Value \\
\midrule
Records & $9{,}218$ \\
Columns & $28$ \\
Unique CVEs & $9{,}218$ \\
Target in soft-label runs & EPSS score \\
Main TPG text & CVE description \\
Summary columns available & all-sources, GitHub-URLs, CVSS-metrics \\
All-sources summary non-empty & $7{,}606$ ($82.5\%$) \\
GitHub-URLs summary non-empty & $2{,}246$ ($24.4\%$) \\
CVSS-metrics summary non-empty & $9{,}218$ ($100\%$) \\
Records with public code & $2{,}339$ \\
\bottomrule
\end{tabular}
\caption{GPT summary dataset.}
\end{table}

\subsection{Gemma Summary Dataset}

Gemma has the same row count, the same CVE coverage and the same
$28$-column schema as GPT. The difference is the summary text: this
dataset's summaries are generated by Gemma.

\begin{table}[h]
\centering
\small
\begin{tabular}{ll}
\toprule
Property & Value \\
\midrule
Records & $9{,}218$ \\
Columns & $28$ \\
Unique CVEs & $9{,}218$ \\
Target in soft-label runs & EPSS score \\
Main TPG text & CVE description \\
All-sources summary non-empty & $7{,}636$ ($82.8\%$) \\
GitHub-URLs summary non-empty & $2{,}381$ ($25.8\%$) \\
CVSS-metrics summary non-empty & $9{,}218$ ($100\%$) \\
Records with public code & $2{,}339$ \\
\bottomrule
\end{tabular}
\caption{Gemma summary dataset.}
\end{table}

Gemma is the controlled counterpart to GPT: the structured fields are
identical, only the LLM that produced the summaries is different. In
baseline runs only the description is used. In summary runs the
relevant summary column is mapped to the unified \emph{LLM summary}
field before the TPG is built.

\subsection{Mistral Summary Dataset (New)}

Mistral was added in this round. Same row count, same CVE coverage,
same $28$-column schema as GPT and Gemma, same three summary columns.

\begin{table}[h]
\centering
\small
\begin{tabular}{ll}
\toprule
Property & Value \\
\midrule
Records & $9{,}218$ \\
Columns & $28$ \\
Unique CVEs & $9{,}218$ \\
Target in soft-label runs & EPSS score \\
Main TPG text & CVE description \\
All-sources summary non-empty & $7{,}317$ ($79.4\%$) \\
GitHub-URLs summary non-empty & $2{,}388$ ($25.9\%$) \\
CVSS-metrics summary non-empty & $9{,}016$ ($97.8\%$) \\
Records with public code & $2{,}339$ \\
\bottomrule
\end{tabular}
\caption{Mistral summary dataset.}
\end{table}

\subsection{DeepSeek Summary Dataset}

DeepSeek has the same row count and CVE coverage but uses the older
$27$-column schema: the date / time / text / source-count /
code-available / delta-days fields take the place of the newer
date-posted / time-posted / social-media-post / occurrence-count /
GitHub-links-with-code-available / days-since-git-source fields. It
also has only two of the three summary columns: the all-sources and
GitHub-URLs summaries are present, the CVSS-metrics summary is not.
The dataset adapter handles either schema, so DeepSeek still works in
the pipeline; the \textsc{S\_cvss} variant is the one that cannot run
on it.

\begin{table}[h]
\centering
\small
\begin{tabular}{ll}
\toprule
Property & Value \\
\midrule
Records & $9{,}218$ \\
Columns & $27$ (older schema) \\
Unique CVEs & $9{,}218$ \\
Target in soft-label runs & EPSS score \\
Main TPG text & CVE description \\
All-sources summary non-empty & $7{,}636$ ($82.8\%$) \\
GitHub-URLs summary non-empty & $2{,}381$ ($25.8\%$) \\
CVSS-metrics summary non-empty & column not present \\
Records with public code & $2{,}339$ \\
\bottomrule
\end{tabular}
\caption{DeepSeek summary dataset. Older schema, no CVSS-metrics
summary.}
\end{table}

\subsection{Megavul Datasets (New)}
\label{sec:megavul}

The Megavul datasets come from the Megavul branch of the same upstream
repository. There are three datasets: Mega-GPT, Mega-Mistral and
Mega-Gemma, each with the same $9{,}486$ rows covering $9{,}240$
unique CVEs. Rows are not one-per-CVE: a single CVE can appear in
several rows when the upstream Megavul corpus links it to more than
one vulnerable commit. The pipeline deduplicates by CVE-ID when
building the labelled records.

The schema is different from the social-media datasets. There is no
social-media post field, no observation date, no source-platform
column. Instead, every row carries a commit URL, the pre-patch and
post-patch function code, and CVSS information in two flavours: the
original NVD CVSS, plus an LLM-augmented v3 estimate when the original
is v2-only. The summary columns are also different: the three LLM
summaries are a commit-URL summary, a code summary (an LLM summary of
the vulnerable function), and a CVSS-metrics summary. There is no
all-sources or GitHub-URLs summary.

\begin{table}[h]
\centering
\small
\begin{tabular}{lrrr}
\toprule
Property & Mega-GPT & Mega-Mistral & Mega-Gemma \\
\midrule
Raw rows                       & $9{,}486$ & $9{,}486$ & $9{,}486$ \\
Unique CVEs                    & $9{,}240$ & $9{,}240$ & $9{,}240$ \\
Columns                        & $29$      & $29$      & $29$ \\
Source language: C/C++ rows    & $8{,}593$ & $8{,}593$ & $8{,}593$ \\
Source language: Java rows     & $893$     & $893$     & $893$ \\
Commit-URL summary non-empty   & $9{,}469$ ($99.8\%$) & $9{,}486$ ($100\%$) & $9{,}486$ ($100\%$) \\
Code summary non-empty         & $9{,}256$ ($97.6\%$) & $9{,}486$ ($100\%$) & $9{,}486$ ($100\%$) \\
CVSS-metrics summary non-empty & $9{,}486$ ($100\%$)  & $9{,}305$ ($98.1\%$) & $9{,}486$ ($100\%$) \\
\bottomrule
\end{tabular}
\caption{Megavul dataset sizes and summary-column coverage. The three
datasets share the same row set; only the LLM that produced the
summaries differs.}
\end{table}

\subsubsection{Megavul Label Inventory}

The Megavul datasets carry several columns that could in principle act
as a training label. Their coverage and class balance are summarised
below (percentages are over the raw $9{,}486$ rows).

\begin{table}[h]
\centering
\small
\begin{tabular}{lrrl}
\toprule
Candidate label column & Non-empty rows & Coverage & Comment \\
\midrule
KEV date-added & $37$ & $0.39\%$ & CISA KEV catalog membership; hard binary signal. \\
CVSS base severity & $9{,}486$ & $100\%$ & 4-class categorical: LOW / MEDIUM / HIGH / CRITICAL. \\
CVSS v3-augmented severity & $3{,}510$ & $37.0\%$ & v3 augmented severity, only for CVEs whose original NVD CVSS was v2. \\
EPSS at publication & varies & partial & EPSS at CVE publication time; often missing. \\
EPSS most-recent & $6{,}694$ & $70.6\%$ & EPSS at the most recent recalculation; the pipeline's primary EPSS source. \\
\bottomrule
\end{tabular}
\caption{Label-candidate columns in the Megavul datasets.}
\end{table}

Two facts dominate the choice of training target.

\begin{itemize}
    \item KEV membership only marks $37$ Megavul rows (about $22$
    unique CVEs once duplicates are collapsed). After a $70/15/15$
    split that leaves roughly three positives in the test set, which
    is too sparse to estimate precision, recall or F1 with any
    confidence. KEV is not a viable training label on Megavul on its
    own.
    \item EPSS coverage is $70.6\%$ when the most-recent EPSS value
    is preferred and the publication-time EPSS is used as a fallback.
    The remaining $29.4\%$ of rows have no EPSS at all and are dropped
    by the dataset adapter.
\end{itemize}

The Megavul soft-label runs use the EPSS score as the continuous
regression target, exactly like the social-media runs. The binary
label used for stratified splitting and for threshold-based metrics
is the same EPSS-thresholded value: positive iff EPSS $\geq 0.10$.

\paragraph{Positive/negative ratio at several thresholds.}
The EPSS distribution is heavily skewed on Megavul; the positive rate
is about seven times sparser than on the social-media datasets. The
numbers below are over the $9{,}486$ rows of the Mega-GPT dataset; the
other two Megavul datasets give the same numbers because EPSS is taken
from the shared NVD/EPSS metadata, not from the LLM summaries.

\begin{table}[h]
\centering
\small
\begin{tabular}{lrrrr}
\toprule
EPSS binarisation threshold & Positives & Negatives & Positive rate & Comment \\
\midrule
$\geq 0.01$ & $1{,}840$ & $7{,}646$ & $19.40\%$ & Very lenient; includes most non-trivial scores. \\
$\geq 0.05$ & $280$     & $9{,}206$ & $2.95\%$  & Moderate. \\
$\geq 0.10$ & $164$     & $9{,}322$ & $1.73\%$  & Pipeline default. \\
$\geq 0.20$ & $124$     & $9{,}362$ & $1.31\%$  & Tight. \\
$\geq 0.50$ & $94$      & $9{,}392$ & $0.99\%$  & Very tight. \\
\bottomrule
\end{tabular}
\caption{Megavul EPSS positive/negative split at several binarisation
thresholds. The pipeline default of $0.10$ yields a $\sim 1.73\%$
positive rate.}
\end{table}

The skew matters for training. With a $\sim 2\%$ positive rate, a
naive classifier that always predicts ``negative'' already reaches
$98\%$ accuracy and the random-baseline PR-AUC equals the positive
rate ($\approx 0.02$). Reported metrics on Megavul should be read
against that baseline, not against the much denser ($\sim 15.5\%$)
social-media baseline.

\paragraph{Why KEV membership is reported but not used as the training label.}
Even with only $\sim 22$ KEV CVEs in Megavul, the column is informative
and is preserved in the converted records. It is useful for
spot-checking the model's high-EPSS predictions against catalog
ground truth: do the KEV CVEs receive high predicted scores? But it
is too sparse to drive optimisation directly. The EPSS-thresholded
binary label, while noisier, has roughly seven times more positives
and is what the soft-label runs report against.

\subsubsection{Megavul Variant Map}

The $15$ Megavul runs in this round use the same five-way ablation
pattern as the $20$ social-media runs. The summary mapping is
different.

\begin{table}[h]
\centering
\small
\begin{tabular}{lll}
\toprule
Variant & Description & LLM summary used \\
\midrule
\textsc{D}        & yes & none \\
\textsc{S\_url}   & no  & commit-URL summary \\
\textsc{S\_code}  & no  & code summary (vulnerable function) \\
\textsc{S\_cvss}  & no  & CVSS-metrics summary \\
\textsc{ALL}      & yes & all three concatenated \\
\bottomrule
\end{tabular}
\caption{Megavul variant map. Fifteen runs total: three datasets
$\times$ five variants. \textsc{S\_url} and \textsc{S\_code} replace
the \textsc{S\_all} and \textsc{S\_git} labels used on the
social-media side, because those columns do not exist in the Megavul
schema.}
\end{table}

\subsection{NVD/KEV Dataset}

The NVD/KEV dataset is built directly from NVD records, the CISA KEV
catalog, FIRST EPSS scores, and the ExploitDB snapshot. Unlike the
LLM-summary corpora, it does not come from a third-party CSV; the
data-collection step assembles it.

\begin{table}[h]
\centering
\small
\begin{tabular}{lrrr}
\toprule
File family & Records & KEV positives & Positive rate \\
\midrule
Full labelled records & $132{,}322$ & $803$ & $0.61\%$ \\
Balanced labelled records (v1) & $4{,}015$ & $803$ & $20.00\%$ \\
Balanced labelled records (v2) & $4{,}015$ & $803$ & $20.00\%$ \\
\bottomrule
\end{tabular}
\caption{NVD/KEV dataset sizes.}
\end{table}

The full file covers CVEs published from $2020$ through $2024$. The
balanced files keep all $803$ KEV positives and sample $3{,}212$
non-KEV negatives. The earlier multiview-hybrid runs used the second
balanced file (the v2 file).

\paragraph{KEV catalog evolution.}
The $803$ figure is fixed at the time the labels were generated; the
local CISA KEV snapshot at that time held $1{,}554$ entries. The
current authoritative CISA KEV catalog has $1{,}590$ CVEs
(catalogVersion 2026.05.08, released 2026-05-08). Re-binding the
labels to the current KEV would require rerunning the data-collection
step and rebuilding the affected experiments; we have not done so for
this round. The KEV cross-reference reported alongside model
predictions in the experiment results document uses the current
$1{,}590$-CVE catalog so that the model's high-confidence non-KEV
predictions can be checked against the up-to-date catalog.

\begin{table}[h]
\centering
\small
\begin{tabular}{p{0.32\textwidth}p{0.56\textwidth}}
\toprule
Source & Purpose \\
\midrule
NVD records dump & Raw NVD records, flattened to description, publication date, CVSS and CWE. \\
CISA KEV catalog snapshot & The binary label source. \\
FIRST EPSS dump & EPSS scores and (in the bulk path) percentiles. \\
ExploitDB snapshot & Public exploit and PoC metadata used by the v2 balanced file. \\
Merged labelled records & Final per-CVE training records consumed by the dataset module. \\
\bottomrule
\end{tabular}
\caption{Sources used to build the NVD/KEV dataset.}
\end{table}

\subsubsection{NVD/KEV Label Generation}

The binary label is dataset-dependent. For the NVD/KEV file it is
generated from CISA KEV membership:

\[
\texttt{binary\_label}_{\text{KEV}} =
\begin{cases}
1, & \text{if CVE-ID is in CISA KEV} \\
0, & \text{otherwise.}
\end{cases}
\]

The same membership signal is also stored as a separate \emph{in-KEV}
flag. For the records we have, the binary label and the in-KEV flag
agree on every record, so there is no label conflict.

For the LLM-summary datasets the same binary-label field exists but
is generated differently. Those datasets carry no KEV signal, so the
adapter thresholds the EPSS score at $0.10$:

\[
\texttt{binary\_label}_{\text{EPSS}} =
\begin{cases}
1, & \text{if EPSS} \geq 0.10 \\
0, & \text{otherwise.}
\end{cases}
\]

The text-source ablations do not actually train against this binary
label. They use the soft-label mode with the raw EPSS score as a
continuous target in $[0, 1]$. The thresholded binary label is only
used for stratified splitting and for reporting precision, recall and
F1 at a $0.5$ decision threshold. The two definitions of binary label
should not be conflated when comparing NVD/KEV runs against
LLM-summary runs.

\subsubsection{NVD/KEV Feature Fields}

\small
\begin{longtable}{p{0.24\textwidth}p{0.18\textwidth}p{0.46\textwidth}}
\caption{Fields in the NVD/KEV labelled dataset.}\\
\toprule
Field & Source & Meaning \\
\midrule
\endfirsthead
\toprule
Field & Source & Meaning \\
\midrule
\endhead
\bottomrule
\endfoot

CVE identifier & NVD & CVE-ID, used as the sample key. \\
Description & NVD & Main vulnerability text used to build the TPG. \\
Published timestamp & NVD & CVE publication timestamp; used for the tabular age feature. \\
CVSS v3 score & NVD & Normalised in the tabular branch. Missing for $0.25\%$ of the full file. \\
CVSS v3 vector & NVD & Vector string parsed into $22$ one-hot tabular dimensions. \\
CWE list & NVD & CWE identifiers, encoded as top-K multi-hot features in hybrid mode. Empty for $11.40\%$ of the full file. \\
References & NVD & Reference URLs. Present in the collector schema, but absent from the saved files. \\
Binary label & CISA KEV & Target for binary-label runs. $1$ means the CVE is in KEV. \\
In-KEV flag & CISA KEV & Boolean copy of the same KEV membership signal. \\
EPSS score & FIRST EPSS & EPSS probability. In binary KEV experiments this is an optional input feature, not the label. \\
EPSS percentile & FIRST EPSS & EPSS percentile. Present in the v2 balanced file; absent from the older full file. \\
High-EPSS flag & derived from EPSS & Convenience flag using a $0.5$ EPSS cutoff. \\
Public-exploit flag & ExploitDB & Whether public PoC or exploit code was found. Present in the v2 balanced file. \\
Number of exploits & ExploitDB & Number of mapped public exploit records. v2 only. \\
Verified-exploit flag & ExploitDB & Whether at least one mapped ExploitDB entry is verified. v2 only. \\
Exploit types & ExploitDB & Exploit category list (remote, webapps, local, DoS, etc.). v2 only. \\
Earliest exploit date & ExploitDB & Earliest public exploit date when available. v2 only. \\

\end{longtable}
\normalsize

\subsubsection{Pipeline Use}

The NVD/KEV dataset is the default training input when no
\verb|--source-csv| is provided. Two common modes:

\begin{itemize}
    \item Without \verb|--skip-collect|, the pipeline fetches NVD,
    KEV, EPSS and ExploitDB data, then writes the merged labelled
    records.
    \item With \verb|--skip-collect|, the pipeline loads existing
    labelled records, either from the default location or from a path
    passed via \verb|--labeled-file|.
\end{itemize}

The LLM-summary datasets are read instead via the \verb|--source-csv|
flag. The dataset adapter maps the relevant raw columns to the
labelled-records format, picks the chosen summary column based on
\verb|--summary-source|, and writes per-run labelled records.

\subsubsection{Robustness Check}

The saved NVD/KEV records pass the basic consistency checks:

\begin{itemize}
    \item every positive label in the full file is present in the KEV
    catalog snapshot;
    \item every KEV CVE present in the full NVD file is labelled
    positive;
    \item the binary label and the in-KEV flag agree on every record;
    \item the balanced files are subsets of the full labelled file;
    \item the v2 balanced file contains all $803$ positives;
    \item descriptions are non-empty after rejected CVEs are removed;
    \item EPSS scores are in the valid $[0, 1]$ range.
\end{itemize}

A few caveats remain:

\begin{itemize}
    \item The balanced files are saved artefacts. The exact
    negative-sampling step is not part of the main pipeline, so these
    files should be treated as fixed unless a new reproducible sampler
    is added.
    \item KEV is conservative. A label of $0$ means ``not in CISA
    KEV'', not ``never exploited''.
    \item KEV labels reflect the catalog snapshot at collection time.
    For temporal experiments this matters because a CVE can be added
    to KEV after its publication date.
    \item The dataset cache hashes the labelled records and
    invalidates processed graph caches when the source changes. This
    matters because older runs reused the same cache root for
    different labelled inputs.
\end{itemize}

\section{Raw Dataset Columns}

The schema below covers the new format used by GPT, Gemma and Mistral.
Where DeepSeek differs (older naming), the difference is noted in the
``Pipeline use'' column.

\small
\begin{longtable}{p{0.20\textwidth}p{0.14\textwidth}p{0.32\textwidth}p{0.28\textwidth}}
\caption{Raw dataset columns and how they are used.}\\
\toprule
Field & Type & Meaning & Pipeline use \\
\midrule
\endfirsthead
\toprule
Field & Type & Meaning & Pipeline use \\
\midrule
\endhead
\bottomrule
\endfoot

CVE identifier & string & The CVE-ID. & Used as the record key, not as a predictive feature. \\

Source platform & categorical & Collection source for the row (Telegram, Reddit, Hacker News, ExploitDB, Mastodon, BleepingComputer, etc.). & Stored as metadata. Not encoded by the model. \\

Date posted & date string & Observation date from the source collection. (DeepSeek uses a renamed date field.) & Combined with the time field to form the publication timestamp. Used by the tabular age feature. \\

Time posted & time string & Time component of the observation. (DeepSeek uses a renamed time field.) & Combined with the date field to form the publication timestamp. \\

Social-media post & string & Original social or source text mentioning the CVE. (DeepSeek uses a renamed text field.) & Not used by the TPG pipeline. The TPG text comes from the description and optionally from an LLM summary. \\

EPSS score & float & FIRST EPSS exploitation probability. & Used as the soft label in soft-label mode. Removed from the input by \verb|--no-epss-feature| to prevent target leakage. \\

EPSS status & categorical & Marks whether the EPSS value is original or enriched. & Not used by the model. Useful for dataset auditing. \\

Description & string & NVD vulnerability description. & The main TPG text. Excluded only by \verb|--summary-only-tpg|. \\

CVSS version & numeric & CVSS version (3.0, 3.1). & Not encoded directly. The adapter rebuilds the CVSS vector from the score and component columns. \\

CVSS score & float & CVSS base score, 0 to 10. & Normalised in the tabular branch. \\

Attack Vector & categorical & CVSS attack vector (NETWORK, LOCAL, ADJACENT\_NETWORK, PHYSICAL). & Rebuilt into the CVSS vector string and one-hot encoded. \\

Attack Complexity & categorical & CVSS attack complexity (LOW, HIGH). & Rebuilt and one-hot encoded. \\

Privileges Required & categorical & Privileges required (NONE, LOW, HIGH). & Rebuilt and one-hot encoded. \\

User Interaction & categorical & User interaction needed (NONE, REQUIRED). & Rebuilt and one-hot encoded. \\

Scope & categorical & Whether exploitation changes scope (UNCHANGED, CHANGED). & Rebuilt and one-hot encoded. \\

Confidentiality impact & categorical & Confidentiality impact (NONE, LOW, HIGH). & Rebuilt and one-hot encoded. \\

Integrity impact & categorical & Integrity impact (NONE, LOW, HIGH). & Rebuilt and one-hot encoded. \\

Availability impact & categorical & Availability impact (NONE, LOW, HIGH). & Rebuilt and one-hot encoded. \\

Occurrence count & integer & Number of source or social records associated with the CVE. (DeepSeek uses a renamed source-count field.) & Mapped to the social-source count and the exploit count. Used as a fallback signal-richness feature. \\

Sources available & boolean & Whether source links or source text were available. & Not encoded by the model. \\

Code-available flag & boolean/integer & Whether public exploit code or PoC is available on GitHub. (DeepSeek uses a renamed code-available field.) & Mapped to the public-exploit flag. Used as a binary tabular feature. \\

Days since latest GitHub source & float & Days since the most recent linked GitHub source. (DeepSeek uses a renamed delta-days field.) & Not currently used by the model. \\

Days since oldest GitHub source & float & Days since the oldest linked GitHub source. (DeepSeek uses a renamed delta-days field.) & Not currently used by the model. \\

Source links & string/list & URLs or source links used as evidence for the CVE row. & Not encoded directly. Upstream evidence for the summary columns. \\

All-sources summary & string & LLM summary built from every available source link. & Selected by \verb|--summary-source all_sources|. Combined into the LLM summary by \verb|--summary-source combined|. \\

GitHub URLs & string/list & Raw GitHub URLs for the CVE. (Not present in DeepSeek.) & Upstream evidence for the GitHub-URLs summary. Not encoded directly. \\

GitHub-URLs summary & string & LLM summary built from GitHub URLs only. & Selected by \verb|--summary-source github_urls|. Combined by \verb|--summary-source combined|. \\

CVSS-metrics summary & string & LLM summary written from the CVSS metrics in templated form. (Not present in DeepSeek.) & Selected by \verb|--summary-source cvss_metrics|. Combined by \verb|--summary-source combined|. \\

\end{longtable}
\normalsize

\section{Converted Fields in the Labelled Records}

\small
\begin{longtable}{p{0.22\textwidth}p{0.32\textwidth}p{0.24\textwidth}p{0.14\textwidth}}
\caption{Fields written by the dataset adapter and consumed by the
TPG dataset.}\\
\toprule
Converted field & Source field(s) & Meaning & Used by model \\
\midrule
\endfirsthead
\toprule
Converted field & Source field(s) & Meaning & Used by model \\
\midrule
\endhead
\bottomrule
\endfoot

CVE identifier & raw CVE field & Sample identifier. & Metadata only. \\
Description & raw description & NVD text used to build the TPG. & Yes, except in \verb|--summary-only-tpg| mode. \\
Publication timestamp & date + time fields & Timestamp proxy. & Yes; used for the tabular age feature. \\
CVSS v3 score & raw CVSS score & CVSS v3 score. & Yes, when the hybrid head is enabled. \\
CVSS v3 vector & CVSS component columns & Reconstructed CVSS vector string. & Yes; parsed into the $22$ one-hot dims. \\
CWE list & not provided & CWE list. Empty in these datasets. & Encoded as zeros (CWE features). \\
References & not provided & NVD reference list. Empty. & Falls back to the source count. \\
Binary label & raw EPSS thresholded at $0.10$ & Binary high-risk label. & Used for stratified splits and threshold-$0.5$ metrics, not as the training target. \\
EPSS score & raw EPSS & Soft EPSS target. & Label in soft-label runs. \\
EPSS percentile & default $0.0$ & Percentile is unavailable here. & Removed by \verb|--no-epss-feature|. \\
High-EPSS flag & raw EPSS $\geq 0.10$ & High-EPSS convenience flag. & Not encoded directly. \\
In-KEV flag & default false & KEV status unavailable in these datasets. & Not used here. \\
Public-exploit flag & raw code-available field & Whether public exploit or PoC code is available. & Yes; binary tabular feature. \\
Exploit count & raw occurrence count & Count proxy. & Yes; log-normalised. \\
Social-source count & raw occurrence count & Same count, kept as a separate field. & Yes; fallback for the references count. \\
Code-available flag & raw code-available field & Boolean, kept as a separate field. & Used as a fallback. \\
LLM summary & one of the three summary columns, or all three joined & Unified summary text. & Yes, when summary-related TPG flags are set. \\
Source platform & raw source field & Original source platform. & Metadata only. \\

\end{longtable}
\normalsize

\section{TPG Features}

The graph branch does not consume raw columns directly. It converts
text into a graph whose nodes represent documents, sentences, tokens,
entities, predicates, arguments, phrases, clauses and security-specific
objects. Edges represent containment, sequence, syntax, semantic
relations, coreference, discourse, and (when enabled) security
relations.

\begin{table}[h]
\centering
\small
\begin{tabular}{p{0.22\textwidth}p{0.28\textwidth}p{0.38\textwidth}}
\toprule
Graph feature & Created from & Dimensions / values \\
\midrule
Node-type one-hot & TPG schema (13 base types) & 13-dim indicator over \textsc{document}, \textsc{paragraph}, \textsc{sentence}, \textsc{token}, \textsc{entity}, \textsc{predicate}, \textsc{argument}, \textsc{concept}, \textsc{noun\_phrase}, \textsc{verb\_phrase}, \textsc{clause}, \textsc{mention}, \textsc{topic}. \\
Text embedding & SecBERT (jackaduma/SecBERT, frozen) & 768-dim contextual embedding per node. Nodes without text get a zero vector. \\
Source one-hot (optional) & \verb|--add-source-labels| & 3-dim indicator: description / summary / mixed. \\
\verb|edge_index| & TPG edges & $2 \times M$ integer tensor with directed connectivity. Used by PyG message passing. \\
\verb|edge_type| & Edge relations & $M$-vector of integer relation indices. $13$ base types ($0\ldots 12$); $23$ total when \verb|--include-security-edges| adds the SEC relations ($13\ldots 22$). \\
\verb|edge_attr| & \verb|edge_type| & One-hot edge-type vector ($13$- or $23$-dim). \\
Security edges & \verb|--include-security-edges| & Adds the SEC relations \textsc{affects}, \textsc{has\_version}, \textsc{located\_in}, \textsc{classified\_as}, \textsc{exploited\_by}, \textsc{causes}, \textsc{mitigated\_by}, \textsc{uses\_function}, \textsc{threatens}, \textsc{has\_severity} ($10$ types). \\
Source labels & summary / two-view flags & Marks nodes and edges as description, summary or mixed. \\
Summary pooling node & \verb|--summary-pooling-node| & Adds one synthetic node linked to the document via \textsc{contains}, with a pooled summary embedding. \\
\bottomrule
\end{tabular}
\caption{Main features used by the TPG branch. The base node feature
dimension is $d_{\text{in}} = 13 + 768 = 781$, or $784$ with source
labels appended.}
\end{table}

The edge features in this table are graph features. They live on TPG
edges and are consumed by edge-aware GNN backbones during message
passing. They are not part of the tabular feature vector. The
multi-view TPG formula document covers the per-view encoder, the
node-conditioned fusion, and the pooling step that produce the graph
embedding.

\section{Tabular Features}

When \verb|--hybrid| is set, the graph embedding is concatenated with
a tabular feature vector $\mathbf{f}_{\text{tab}}$. After the
leakage-fix pass, the default configuration is $55$-dim: the EPSS
score and percentile are removed from the input because the soft
target is EPSS itself. The $57$-dim variant (with EPSS as input) is no
longer the default and is kept only for the NVD/KEV binary
experiments, where the target is KEV membership rather than EPSS.

The tabular vector contains CVE-level metadata: CVSS, CWE,
vulnerability age, source counts and public-exploit availability. It
does not contain edge types or SEC edge labels; those live in the
graph branch.

\begin{table}[h]
\centering
\small
\begin{tabular}{p{0.28\textwidth}rp{0.42\textwidth}}
\toprule
Feature group & Dim & Meaning \\
\midrule
CVSS score and has-CVSS flag & 2 & Normalised CVSS v3 score in $[0, 1]$ plus a binary flag indicating whether CVSS data exists. \\
CVSS vector one-hot & 22 & One-hot encoding of the eight CVSS v3 metrics (see Table~\ref{tab:cvss22}). \\
CWE multi-hot and ``other'' bucket & 26 & Top-$25$ CWE IDs plus a catch-all bucket. Populated for NVD/KEV; all zero for the LLM-summary datasets (no CWE list there). \\
CWE count & 1 & Normalised count of CWE IDs (capped at $10$, divided by $10$). Populated for NVD/KEV; zero for the LLM-summary datasets. \\
Reference or source count & 1 & Log-normalised $\log(1 + n) / \log(21)$ of NVD reference URLs, or social/source mention count when references are absent. \\
Vulnerability age & 1 & Log-normalised age in days from the publication-date proxy, $\log(1 + \text{age}) / \log(3651)$. \\
EPSS score and percentile (optional) & 2 & Direct EPSS inputs. Excluded by default. Included only in the NVD/KEV binary runs where EPSS is not the label. \\
Public-exploit flag & 1 & Binary indicator that public PoC or exploit code is available. \\
Exploit/source count & 1 & Log-normalised count proxy from the occurrence count, with the social-source count as a fallback. \\
\bottomrule
\end{tabular}
\caption{Tabular features used by the hybrid model. Default total is
$\dim \mathbf{f}_{\text{tab}} = 55$ (rows excluding the optional 2-D
EPSS group). With EPSS included it is $57$.}
\end{table}

\begin{table}[h]
\centering
\small
\begin{tabular}{lcl}
\toprule
CVSS v3.1 metric & Cardinality & Possible values \\
\midrule
Attack Vector (AV) & 4 & NETWORK, ADJACENT\_NETWORK, LOCAL, PHYSICAL \\
Attack Complexity (AC) & 2 & LOW, HIGH \\
Privileges Required (PR) & 3 & NONE, LOW, HIGH \\
User Interaction (UI) & 2 & NONE, REQUIRED \\
Scope (S) & 2 & UNCHANGED, CHANGED \\
Confidentiality (C) & 3 & NONE, LOW, HIGH \\
Integrity (I) & 3 & NONE, LOW, HIGH \\
Availability (A) & 3 & NONE, LOW, HIGH \\
\midrule
Total & 22 & \\
\bottomrule
\end{tabular}
\caption{Decomposition of the 22-dim CVSS one-hot.}
\label{tab:cvss22}
\end{table}

\section{Which Features Are Used by Each Experiment}

The social-media text-source ablation uses five mutually exclusive
variants. The five rows below match the variants reported in the
experiment results.

\begin{table}[h]
\centering
\small
\begin{tabular}{lccccc}
\toprule
Variant & Description & LLM summary used & Tabular & EPSS input & Security edges \\
\midrule
\textsc{D}      & yes & none                                        & yes & no & no \\
\textsc{S\_all} & no  & all-sources summary                         & yes & no & no \\
\textsc{S\_git} & no  & GitHub-URLs summary                         & yes & no & no \\
\textsc{S\_cvss}& no  & CVSS-metrics summary (not on DeepSeek)      & yes & no & no \\
\textsc{ALL}    & yes & all three concatenated                      & yes & no & no \\
\bottomrule
\end{tabular}
\caption{Feature use across the five social-media text-source variants.}
\end{table}

The Megavul runs use the same variant scheme but a different summary
column map; see Section~\ref{sec:megavul} for the
\textsc{S\_url} / \textsc{S\_code} / \textsc{S\_cvss} mapping.
Targets, tabular branch and EPSS handling are identical to the
social-media table above.

The other experiment types use different combinations.

\begin{table}[h]
\centering
\small
\begin{tabular}{lccccc}
\toprule
Experiment & Description & LLM summary & Tabular & EPSS input & Security edges \\
\midrule
NVD/KEV binary rerun         & yes & no  & yes      & yes (not the label) & no \\
NVD/KEV temporal shift       & yes & no  & yes      & yes                 & no \\
NVD/KEV no-EPSS extrapolation & yes & no  & yes      & no                  & no \\
\bottomrule
\end{tabular}
\caption{Feature use across the NVD/KEV experiments.}
\end{table}

\section{Leakage Notes}

The main leakage risk in the LLM-summary experiments is using the EPSS
score as an input feature while also predicting EPSS. That puts the
target directly into the input and inflates the metrics. The
leakage-free EPSS soft-label experiments remove this input with
\verb|--no-epss-feature|.

In the NVD/KEV binary experiments the EPSS score is not the target;
KEV membership is. EPSS is still a strong external exploitation-risk
signal, so results with EPSS included should be read as hybrid
prediction with an external risk signal, not as a pure text-only
result.

A second, more subtle risk is that the CVSS-metrics summary is just
templated CVSS prose. The same CVSS information is already in the
$22$-dim one-hot tabular block, so feeding it again as text adds
redundancy and seems to displace the description signal. The
text-source ablations show this empirically: \textsc{S\_cvss} is the
worst single-source variant on every dataset.

Other fields are not labels, but some are strong proxies. CVSS
severity, public-exploit availability, source count and the LLM
summaries all carry information related to exploitation likelihood.
They are useful modelling features, but the reported results should
make clear which of these signals were available to the model.

\end{document}
